% =============================================================================
%  CHAPTER 2 — BACKGROUND & RELATED WORK (~25-35% of thesis)
% =============================================================================
\section{Background and Related Work}
\label{sec:background}

This section provides the theoretical foundations necessary to understand the
system developed in this thesis. We cover large language models, embedding
representations, retrieval-augmented generation, vector databases, and the
application of AI in mathematics education.

% ─────────────────────────────────────────────────────────────────────────────
\subsection{Large Language Models}
\label{subsec:llm-background}

\subsubsection{Transformer Architecture}

The Transformer architecture, introduced by Vaswani et
al.~\cite{vaswani2017attention}, forms the foundation of modern large
language models. Unlike recurrent neural networks, Transformers process
entire input sequences in parallel using a \textit{self-attention} mechanism
that allows each token to attend to all other tokens in the sequence. This
parallelism enables efficient training on large corpora and better modeling
of long-range dependencies.

A Transformer consists of an encoder and a decoder, each built from stacked
layers of multi-head self-attention and feed-forward networks. The
self-attention mechanism computes, for each token, a weighted sum of all
token representations in the sequence, where the weights are determined by
learned query, key, and value projections:

\begin{equation}
\text{Attention}(Q, K, V) = \text{softmax}\left(\frac{QK^T}{\sqrt{d_k}}\right) V
\label{eq:attention}
\end{equation}

where $Q$, $K$, and $V$ are the query, key, and value matrices, and $d_k$
is the dimension of the key vectors. Multi-head attention runs this
computation multiple times in parallel with different learned projections,
allowing the model to attend to information from different representation
subspaces.

\subsubsection{From BERT to Generative Models}

Two main paradigms emerged from the Transformer architecture:

\begin{itemize}
    \item \textbf{Encoder-only models} such as BERT~\cite{devlin2019bert}
    are pre-trained with masked language modeling (predicting randomly masked
    tokens) and are primarily used for understanding tasks: classification,
    named entity recognition, and—crucially for this work—computing text
    embeddings for retrieval.

    \item \textbf{Decoder-only models} such as
    GPT-3~\cite{brown2020gpt3}, LLaMA~\cite{touvron2023llama}, and
    Gemini~\cite{team2023gemini} are trained autoregressively (predicting
    the next token) and excel at text generation. These models form the
    backbone of modern chatbots and are used in this thesis for answer
    generation.
\end{itemize}

The scaling of decoder-only models has led to emergent capabilities:
in-context learning (performing tasks from examples in the prompt without
fine-tuning), chain-of-thought reasoning, and instruction
following~\cite{brown2020gpt3}. However, these models are prone to
\textit{hallucination}—generating plausible but factually incorrect
content—which is particularly dangerous in an educational context where
students trust the system's mathematical reasoning.

\subsubsection{Gemini}

Google's Gemini family~\cite{team2023gemini} represents a class of
multimodal LLMs capable of processing text, images, audio, and video.
In this thesis, we use Gemini models for two distinct purposes:

\begin{enumerate}
    \item \textbf{Document digitization}: Gemini's vision capabilities are
    used to convert scanned mathematical documents (images of handwritten
    and printed exercises) into structured LaTeX, leveraging the model's
    ability to recognize two-dimensional mathematical notation.

    \item \textbf{Answer generation}: Gemini serves as the generation
    backbone in all three system variants (RAG, prompt-only, hybrid),
    receiving retrieved context and producing curriculum-aligned answers.
\end{enumerate}

Using the same model family for both tasks ensures consistency and simplifies
deployment on Google Cloud Platform's Vertex AI infrastructure.


% ─────────────────────────────────────────────────────────────────────────────
\subsection{Text Embeddings and Semantic Similarity}
\label{subsec:embeddings-background}

\subsubsection{Dense Vector Representations}

Text embeddings map variable-length text sequences to fixed-dimensional
dense vectors in a continuous space where semantic similarity corresponds to
geometric proximity. Early approaches derived embeddings from word-level
models (Word2Vec, GloVe), but modern methods use full-sequence encoders
based on Transformers.

Sentence-BERT~\cite{reimers2019sentence_bert} demonstrated that fine-tuning
BERT with a siamese architecture on natural language inference data produces
embeddings where cosine similarity reliably reflects semantic relatedness.
This approach made it practical to use neural embeddings for large-scale
retrieval, as embedding computation is a single forward pass per document
(offline) and per query (online).

\subsubsection{BGE-M3: Multilingual Multi-Granularity Embeddings}

For this thesis, we selected the BGE-M3 model~\cite{chen2024bge_m3}, which
offers three key properties:

\begin{enumerate}
    \item \textbf{Multi-lingual}: BGE-M3 supports over 100 languages,
    including French and Arabic—both essential for Tunisian Bac materials,
    which are written in French with occasional Arabic mathematical
    terminology.

    \item \textbf{Multi-granularity}: The model supports dense retrieval,
    sparse retrieval (lexical matching), and ColBERT-style late interaction,
    enabling flexible retrieval strategies.

    \item \textbf{Long context}: With a maximum sequence length of 8192
    tokens, BGE-M3 can embed entire exercise-correction pairs without
    truncation, preserving the logical structure of mathematical documents.
\end{enumerate}

The choice of BGE-M3 over alternatives such as OpenAI's text-embedding-ada
or Cohere's multilingual embeddings was motivated by the need for a
self-hosted, open-source model that can run locally without external API
calls, ensuring reproducibility and data sovereignty.

\subsubsection{Challenges of Mathematical Embeddings}

Embedding mathematical content poses specific challenges not encountered in
general-purpose text retrieval. Mathematical formulas encode meaning through
spatial structure (fractions, subscripts, exponents) that is lost in linear
text serialization. Prior work by Mansouri et
al.~\cite{mansouri2019tangent} and Pfahler et
al.~\cite{pfahler2020math_embeddings} has explored specialized mathematical
embeddings, but these require custom architectures and are not directly
applicable to the mixed text-and-formula content of Bac exercises.

Our approach sidesteps this limitation by digitizing mathematical content
into LaTeX notation before embedding. LaTeX provides a standardized linear
serialization of mathematical structures (e.g., \texttt{\textbackslash frac\{a\}\{b\}} for
$\frac{a}{b}$) that general-purpose embedding models can process. While
this may not capture all mathematical semantics, the combination of LaTeX
text with rich metadata filtering (chapter, exercise type) compensates for
any loss in pure embedding similarity.


% ─────────────────────────────────────────────────────────────────────────────
\subsection{Retrieval-Augmented Generation}
\label{subsec:rag-background}

\subsubsection{Motivation and Core Idea}

Retrieval-Augmented Generation (RAG) was introduced by Lewis et
al.~\cite{lewis2020rag} to address a fundamental limitation of parametric
language models: their knowledge is frozen at training time and cannot be
updated without retraining. RAG augments a generative model with a
non-parametric retrieval component that fetches relevant documents from an
external knowledge base at inference time, conditioning the generation on
up-to-date, domain-specific information.

The basic RAG pipeline consists of three stages:
\begin{enumerate}
    \item \textbf{Indexing}: Documents are split into chunks, embedded into
    dense vectors, and stored in a vector database.
    \item \textbf{Retrieval}: Given a user query, the system retrieves the
    $k$ most similar document chunks by comparing query and document
    embeddings.
    \item \textbf{Generation}: The retrieved chunks are concatenated with
    the query and fed to a language model, which generates an answer
    grounded in the retrieved evidence.
\end{enumerate}

\subsubsection{Retrieval Methods}

Two main families of retrieval methods exist:

\begin{itemize}
    \item \textbf{Sparse retrieval} methods such as
    BM25~\cite{robertson2009bm25} use term-frequency statistics to match
    queries to documents. They are efficient, interpretable, and robust to
    domain shift, but cannot capture semantic similarity beyond lexical
    overlap.

    \item \textbf{Dense retrieval} methods such as
    DPR~\cite{karpukhin2020dpr} encode queries and documents into dense
    vectors using neural networks and retrieve by nearest-neighbor search.
    Dense methods capture semantic meaning but require training data and
    may fail on out-of-distribution queries.
\end{itemize}

Modern RAG systems often combine both approaches in hybrid retrieval
pipelines. In this thesis, we use dense retrieval with BGE-M3 embeddings,
augmented by metadata filtering that acts as a structured sparse signal
(filtering by chapter, document type, etc.).

\subsubsection{Chunking Strategies}

How documents are split into chunks significantly affects retrieval quality.
Common approaches include fixed-size chunking (split every $n$ characters),
sentence-based splitting, and semantic chunking (split at topic
boundaries)~\cite{zhao2024financial_chunking, chen2024mog}.

Zhao et al.~\cite{zhao2024financial_chunking} demonstrated that
structure-aware chunking—respecting the natural boundaries of document
elements—leads to better retrieval than arbitrary character-level splitting.
Chen et al.~\cite{chen2024mog} further showed that mixing different chunk
granularities can improve retrieval by capturing both fine-grained details
and broader context.

In our system, we adopt an \textit{adaptive chunking} strategy: smaller
chunks (1500 characters) for corrections to preserve individual exercise
boundaries, and larger chunks (3000 characters) for course material to keep
theorem statements and their proofs together.

\subsubsection{Advanced RAG Architectures}

Since the original RAG paper, numerous extensions have been
proposed~\cite{gao2024rag_survey}:

\begin{itemize}
    \item \textbf{RETRO}~\cite{borgeaud2022retro} retrieves from trillions
    of tokens during pre-training, showing that retrieval can improve model
    quality beyond what scaling alone achieves.

    \item \textbf{Atlas}~\cite{izacard2022atlas} jointly trains the
    retriever and generator, demonstrating that few-shot learning with
    retrieval can match or exceed much larger models trained without
    retrieval.

    \item \textbf{Self-RAG} and similar methods add self-reflection
    mechanisms where the model decides when retrieval is needed and
    evaluates the relevance of retrieved documents.
\end{itemize}

Our system contributes to this landscape by introducing a
\textit{confidence-based routing} mechanism in the hybrid engine: rather
than always retrieving or never retrieving, the system evaluates retrieval
quality and dynamically selects between RAG-grounded and prompt-only
generation.


% ─────────────────────────────────────────────────────────────────────────────
\subsection{Vector Databases}
\label{subsec:vector-db-background}

Vector databases are specialized storage systems optimized for similarity
search over high-dimensional vectors. Unlike traditional databases that
support exact matching, vector databases use approximate nearest neighbor
(ANN) algorithms to efficiently find the most similar vectors to a query.

Several vector database solutions exist, including
FAISS~\cite{johnson2019faiss} (a library for similarity search developed by
Meta), Pinecone (a managed cloud service), Weaviate, and
ChromaDB~\cite{chromadb2024}. We selected ChromaDB for this thesis based on
several criteria:

\begin{itemize}
    \item \textbf{Local persistence}: ChromaDB supports fully local storage,
    avoiding dependence on external services and ensuring reproducibility.
    \item \textbf{Metadata filtering}: ChromaDB allows attaching arbitrary
    metadata to documents and filtering retrieval results by metadata fields,
    which we use extensively for chapter-based and document-type filtering.
    \item \textbf{Python-native API}: ChromaDB provides a clean Python API
    that integrates naturally with the rest of our pipeline.
    \item \textbf{Custom embedding functions}: ChromaDB allows plugging in
    custom embedding models, which we use to integrate BGE-M3.
\end{itemize}


% ─────────────────────────────────────────────────────────────────────────────
\subsection{OCR and Mathematical Document Processing}
\label{subsec:ocr-background}

Converting scanned mathematical documents into machine-readable text is a
long-standing challenge in document analysis~\cite{zanibbi2016math_retrieval}.
Standard OCR systems designed for linear text fail on mathematical notation
because formulas are inherently two-dimensional: fractions stack vertically,
subscripts and superscripts require spatial interpretation, and integral
bounds span multiple lines.

Traditional approaches use specialized math-OCR systems such as Mathpix,
which combines neural network-based symbol recognition with grammar-based
expression parsing~\cite{mathpix2024}. More recently, multimodal LLMs have
demonstrated strong mathematical OCR capabilities, as evaluated by
Wang et al.~\cite{wang2023ocrbench} and Mahdavi et
al.~\cite{mahdavi2019icdar_math}.

In this thesis, we use Google Gemini's vision capabilities for digitization
rather than a dedicated math-OCR tool. This choice is motivated by the
observation that Gemini can simultaneously handle the mixed French text,
mathematical formulas, and handwritten annotations present in Tunisian Bac
materials—a combination that would require multiple specialized tools in a
traditional pipeline.


% ─────────────────────────────────────────────────────────────────────────────
\subsection{AI in Mathematics Education}
\label{subsec:ai-education}

The application of AI to mathematics education has gained significant
attention with the advent of capable LLMs~\cite{kasneci2023chatgpt_education}.
Several categories of tools have emerged:

\begin{itemize}
    \item \textbf{General-purpose chatbots} (ChatGPT, Gemini, Claude) can
    solve mathematical problems but lack curriculum alignment, may use
    methods outside a specific educational program, and do not adapt to
    regional pedagogical conventions.

    \item \textbf{Math-specific tools} (Wolfram Alpha, Photomath) excel at
    computation and step-by-step solutions but operate independently of any
    curriculum context.

    \item \textbf{Curriculum-aligned tutoring systems} (Khan Academy's
    Khanmigo, various national platforms) attempt to ground AI assistance
    within specific educational frameworks but are primarily developed for
    English-speaking markets.
\end{itemize}

To our knowledge, no prior work has addressed the specific combination of:
(1)~the Tunisian Bac mathematics curriculum, (2)~the French-language
correction style conventions, and (3)~a RAG-based approach grounded in
actual examination materials. This thesis fills this gap by developing a
domain-specific RAG system that operates within these constraints.


% ─────────────────────────────────────────────────────────────────────────────
\subsection{Summary}
\label{subsec:background-summary}

Table~\ref{tab:background-summary} summarizes the key technologies used in
this thesis and their roles.

\begin{table}[H]
\centering
\caption{Summary of technologies and their roles in the system.}
\label{tab:background-summary}
\begin{tabularx}{\textwidth}{lX}
\toprule
\textbf{Technology} & \textbf{Role in This Thesis} \\
\midrule
Transformer / Self-Attention & Foundation architecture for both embedding and generation models \\
BGE-M3 Embeddings & Multilingual dense embeddings for French/Arabic math content \\
RAG Pipeline & Core architecture: retrieve relevant docs, then generate grounded answers \\
ChromaDB & Local vector database with metadata filtering \\
Gemini (Multimodal LLM) & Document digitization (vision) and answer generation (text) \\
Adaptive Chunking & Domain-specific document splitting strategy \\
\bottomrule
\end{tabularx}
\end{table}
